\documentclass[11pt, aspectratio = 169]{beamer}

\usetheme[progressbar=frametitle]{metropolis}
\usepackage{appendixnumberbeamer}
\usepackage{booktabs}
\usepackage[scale=2]{ccicons}
\usepackage{pgfplots}
\usepgfplotslibrary{dateplot}
\usepackage{xspace}
\usepackage{hyperref}
\usepackage{verbatim} 
\newcommand{\themename}{\textbf{\textsc{metropolis}}\xspace}
\definecolor{links}{HTML}{1B1B81}
\hypersetup{colorlinks,linkcolor=blue,urlcolor=blue}

\title{Applied Economics}
\subtitle{OLS Variance and Robust Standard Errors}
\author{Beate Thies}
\date{}
\institute{Winter 2024/25}


\setbeamercolor{normal text}{fg=black}
\setbeamercolor{block title}{fg=teal}

\let\svdescription\description
\let\svenddescription\enddescription
\renewenvironment{description}{%
  \svdescription%
  \let\svitem\item%
  \renewcommand\item[1][]{\svitem[\textcolor{\descriptioncolor}{##1}]}%
  }{
  \svenddescription%
  }
\newcommand\descriptioncolor{teal}

\usepackage{Sweave}
\begin{document}
\Sconcordance{concordance:AE_w5_OLS.tex:AE_w5_OLS.Rtex:%
1 36 1 1 0 347 1}


\maketitle

\begin{frame}{Plan for this Lecture}
\begin{enumerate}
    \item Variance of the OLS Estimator
     \begin{itemize}
        \item Special case of Homoskedasticity (MLR.5)
        \item Role of Multicollinearity
    \end{itemize}
    \item Empirical Examples in R (robust standard errors)
    \item Some more functionalities in R (ggplot2, writing functions, loops)
\end{enumerate}
\end{frame}
\section{Variance of the OLS Estimator}

\begin{frame}{Variance of OLS}
    To quantify uncertainty and for statistical inference we need to know/estimate the variance of $\hat{\boldsymbol{\beta}}$.

\begin{align*}
Var(\hat{\boldsymbol{\beta}}|\mathbf{X}) =  \begin{pmatrix}
                                            Var(\hat{\beta_0}|\mathbf{X}) & cov(\hat{\beta_0}, \hat{\beta_1}|\mathbf{X})  & \ldots & cov(\hat{\beta_0}, \hat{\beta_k}|\mathbf{X})   \\
                                            cov(\hat{\beta_0}, \hat{\beta_1}|\mathbf{X})  & Var(\hat{\beta_1}|\mathbf{X})  & \ldots & cov(\hat{\beta_1}, \hat{\beta_k}|\mathbf{X})      \\
                                            \vdots & \vdots   & \ddots & \vdots \\
                                            cov(\hat{\beta_0}, \hat{\beta_k}|\mathbf{X})  &  cov(\hat{\beta_1}, \hat{\beta_k}|\mathbf{X})  & \ldots & Var(\hat{\beta_k}|\mathbf{X}) \\
                                            \end{pmatrix}
\end{align*}

We are especially interested in the elements on the diagonal, which reflect the estimation uncertainty regarding the parameters $\beta_j, j = 0,1,...,k$
\end{frame}

\begin{frame}{Variance of OLS}
  Remember, $\hat{\boldsymbol{\beta}}$ can be written as $\hat{\boldsymbol{\beta}} = \boldsymbol{\beta} + (\mathbf{X}'\mathbf{X})^{-1}\mathbf{X}'\mathbf{u}$
  
Thus,  $\hat{\boldsymbol{\beta}} - \boldsymbol{\beta} = (\mathbf{X}'\mathbf{X})^{-1}\mathbf{X}'\mathbf{u}$

This implies: 
\begin{align*}
Var(\hat{\boldsymbol{\beta}}|\mathbf{X}) &= E[(\hat{\boldsymbol{\beta}}-\boldsymbol{\beta})(\hat{\boldsymbol{\beta}}-\boldsymbol{\beta})'|\mathbf{X}]\\
                            &= E[((\mathbf{X}'\mathbf{X})^{-1}\mathbf{X}'\mathbf{u})((\mathbf{X}'\mathbf{X})^{-1}\mathbf{X}'\mathbf{u})'|\mathbf{X}]\\
                            &= E[(\mathbf{X}'\mathbf{X})^{-1}\mathbf{X}'\mathbf{u}\mathbf{u}'\mathbf{X}(\mathbf{X}'\mathbf{X})^{-1}|\mathbf{X}]\\
                            &= (\mathbf{X}'\mathbf{X})^{-1}\mathbf{X}'E[\mathbf{u}\mathbf{u}'|\mathbf{X}]\mathbf{X}(\mathbf{X}'\mathbf{X})^{-1}\\
\end{align*}
Thus, we need the variance matrix of the error terms $Var(\mathbf{u}|\mathbf{X}) = E[\mathbf{u}\mathbf{u}'|\mathbf{X}]$.
\end{frame}


\begin{frame}{Variance of the error terms}
\begin{align*}
Var(\mathbf{u}|\mathbf{X}) &= E(\mathbf{u}\mathbf{u}'|\mathbf{X}) = \begin{pmatrix}
\sigma_1^2 & 0        & \ldots & 0     \\
0        & \sigma_2^2 & \ldots & 0     \\
\vdots   & \vdots   & \ddots & \vdots \\
0        &  0       & \ldots & \sigma_n^2 \\
\end{pmatrix} = diag(\sigma_1^2, ..., \sigma_n^2) \equiv \Omega
\end{align*}
The elements along the diagonal are the unknown variances of the error terms for all observations $i = 1,...,n.$
Why are the elements outside the diagonal all 0?

\pause

Since $\Omega$ is unknown, we need to estimate it. To do so, we replace $\sigma_i^2$ with the squared residual of observation $i$: \begin{align*}
\widehat{\Omega} = diag(\hat{u}_1^2, \hat{u}_2^2, ..., \hat{u}_n^2)
\end{align*}
\end{frame}

\begin{frame}{Variance Estimator}

Plugging $\widehat{\Omega}$ into the formula for $Var(\hat{\boldsymbol{\beta}}|\mathbf{X})$ yields a consistent estimator for the variance matrix of the OLS estimator:
\begin{align*}
   \widehat{Var}(\hat{\boldsymbol{\beta}}|\mathbf{X}) = (\mathbf{X}'\mathbf{X})^{-1}\mathbf{X}'\hat{\Omega}\mathbf{X}(\mathbf{X}'\mathbf{X})^{-1}
\end{align*}
Diagonal elements of this matrix: estimated variances of the (k+1) parameter estimators. \\
Taking the square roots yields the \color{teal}\textbf{standard errors} $\widehat{SE}(\hat{\beta}_j)$ \color{black} which are relevant for hypothesis testing.

\end{frame}


\begin{frame}{Heteroskedasticity vs. Homoskedasticity}
The variance estimator we just derived is called \color{teal}\textbf{heteroskedasticity-robust}  \color{black}estimator. The corresponding standard errors are called heteroskedasticity-robust standard errors.

Per default, software programs like R report different standard errors that are based on a \textbf{homoskedasticity assumption (MLR.5)}.
\vspace{+.12cm}

\begin{alertblock}{MLR.5 Conditional Homoskedasticity}
\vspace{+.12cm}
The errors $u_i$ have constant conditional variance that does not depend on the values of the regressors:
$var(u_i|\mathbf{X_i}) = \sigma_u^2 \text{ for all } i = 1,...n.$ \\
In matrix notation: $var(\mathbf{u}|\mathbf{X}) = \sigma_u^2\mathbm{I}_n$,
\end{alertblock}

Under this assumption the variance of the OLS estimator simplifies to:
$ Var(\hat{\boldsymbol{\beta}}|\mathbf{X}) = (\mathbf{X}'\mathbf{X})^{-1}\sigma_u^2$
\end{frame}

\begin{frame}{Heteroskedasticity vs. Homoskedasticity}
\alert{MLR.5 is implausible in most economic applications.} \\
Example: Among people with high levels of education, the variance in wages is larger than among people with low levels of education. $\rightarrow$ Unexplained variation in wages ($Y$) depends on education ($X$).  

If MLR.5 is violated, the non-robust standard errors based on the simplified formula are inconsistent. The st. errors and t-statistics reported by standard software are invalid.

The robust standard errors are consistent under both conditional homoskedasticity  and conditional heteroskedasticity. Typically, under MLR.5, they are a bit more conservative than necessary.
\end{frame}

\begin{frame}{Heteroskedasticity-robust standard errors}
Heteroskedasticity is viewed as the norm in most economic applications. Thus, the current empirical literature relies almost exclusively on robust standard errors. 

In R, you can get the robust variance estimator using \textcolor{teal}{\texttt{vcovHC(model, type = "HC1")}} from the package \textcolor{teal}{\texttt{sandwich}}.

To report robust st. errors in the regression output, use \textcolor{teal}{\texttt{coeftest(model, vcov. = vcovHC(model, type = "HC1")}} from the package \textcolor{teal}{\texttt{lmtest}}.
\end{frame}


\begin{frame}{Drivers of estimation uncertainty}
We will assume homoskedasticity for a second, to analyse which factors drive the variance of the OLS estimator.
(Also holds under heteroskedasticity, but is more difficult to show).

Variance of estimator $\hat{\beta_j}$ under homoskedasticity:
\begin{align*}
    Var(\hat{\beta_j}|\mathbf{X}) = \frac{1}{n} \frac{\sigma_u^2}{\sigma_{X_j}^2}\frac{1}{1-R_j^2}
\end{align*}
where $R_j^2$ is the $R^2$ from the regression of $X_j$ on all other regressors $X_l$, $l \neq j$.

  \begin{itemize}
     \item shrinks as sample size \textcolor{teal}{$n$} increases.
     \item shrinks as \textcolor{teal}{sample variation in regressor $X_j$} increases.
     \item increases in the error variance \textcolor{red}{$\sigma_u^2$}.
    \item increases as the \textcolor{red}{correlation between $X_j$ and the other regressors} increases..
 \end{itemize}
\end{frame}


\begin{frame}{Excursion: $R^2$ (goodness of fit)}

$R^2$ is commonly reported in regression output in R, and measures what share of the total variation in the dependent variable is explained by variation in the regressors:
\begin{align*}
 R^2=\frac{SSE}{SST}=1-\frac{SSR}{SST}, \text{ with }
 \end{align*}
\textbf{Total sum of squares} (SST): $\sum_{i=1}^N (y_i - \bar{y})^2$\\
\textbf{Explained sum of squares} (SSE): $\sum_{i=1}^N (\hat{y}_i - \bar{y})^2$\\
\textbf{Residual sum of squares} (SSR): $\sum_{i=1}^N \hat{u}_i^2$


The (adjusted) $R^2$ is sometimes used to compare different models (for the same dependent variable $Y$), but it does not play a major role when your goal is causal analysis.
\end{frame}

\begin{frame}{Multicollinearity and the OLS variance}
$R_j^2$ is the $R^2$ from a regression of $X_j$ on the other regressors. \\
$\rightarrow$ measures multicollinearity (linear correlation) between $X_j$ and the other regressors. 
\begin{itemize}
    \item $0 \leq R_j^2 \leq 1$.
    \item $ R_j^2 =1$ ruled out by MLR.3 (no perfect multicollinearity). 
\end{itemize}

Estimates are less precise (higher variance) if there is strong correlation among regressors.
\begin{itemize}
    \item Intuition: If two regressors are strongly correlated, it is hard to disentangle their individual partial effects on $Y$.
\item $\frac{1}{1-R_j^2}$: Variance Inflation Factor
\end{itemize}
\end{frame}

\section{Getting Robust standard errors in R}

\begin{frame}[fragile]{Non-robust St. Errors}
\small
<<include=FALSE>>=
library(ggplot2)
 wage1 = read.csv("wage1.csv")
@
\footnotesize
<<cache = F, message= F>>=
library(sandwich) # for robust variance estimator
library(lmtest)   # for regression tables

reg_model <- lm(lwage ~ educ + exper + I(exper^2) + female, data = wage1)
coeftest(reg_model) # uses homoskedastic errors! [same as summary(reg_model)]
@
\end{frame}

\begin{frame}[fragile]{Robust St. Errors}

\footnotesize
<<cache = F, message= F>>=
# get the robust variance estimate
var_robust <- vcovHC(reg_model, type = "HC1") 
# regression table based on the robust standard errors
coeftest(reg_model, vcov. = var_robust) 
@
\end{frame}


\section{Creating Graphs with ggplot2}

\begin{frame}{ggplot2 syntax}
\texttt{\color{teal}ggplot(data) + \color{red}<GEOM FUNCTION>(\color{blue}aes(<MAPPINGS>)\color{red})}
    
    
\textbf{Key components}
\begin{itemize}
    \item \color{teal} dataset containing the variables we want to plot 
    \item \color{red} one or more geom functions: the type of object we can observe in the plot, e.g. points, lines, boxplots, or bars.
    \item \color{vlue} aesthetic attributes of the geometric object, e.g. x/y position, color,  or shape. Aesthetic attributes are mapped to the variables in the dataset.
\end{itemize}
\end{frame}



\begin{frame}{ggplot2 syntax}
   
    \begin{itemize}
    \item \textt{geom\_functions: geom\_point, geom\_smooth, geom\_density, geom\_col, ...}
    \item in each geom functions, we have to specify the required aesthetic attributes, e.g. x- and y-variable, color, etc.
    \item to optimize the \alert{appearance of the plot} we can change axes labels (\texttt{xlab("..."), ylab("..."))}, choose a different theme (e.g. \texttt{theme\_light(), theme\_dark(), theme\_minimal()}) or add a title ( \texttt{ggtitle("...")}). These \alert{options are all added} to \textt{ggplot(data) + geom\_function(aes(...))}
\end{itemize}
\end{frame}


\begin{frame}[fragile]{Using Loops in R}  \small
<<cache = F, message= F,out.width='2.2in'>>=
for (i in 1:10) {   
  # loop body: gets executed each time 
  # the value of i changes with each iteration 
}
@ 

The index does not have to be called \texttt{i}.\\

You cannot only loop through \texttt{1:10}, but also any other numbers collected in a vector.
\end{frame}

\begin{frame}[fragile]{Using Loops in R: Examples} \small
<<cache = F, message= F,out.width='2.2in'>>=

for (j in c(1,3,5)) {  # here we call the index j 
    print(j)           # first print the current number
    print(j + 10)      # then add 10 to it and print the result
}
@ 
\end{frame}

\begin{frame}[fragile]{Using Loops in R: Examples}  \small
<<cache = F, message= F,out.width='2.2in'>>=
for (j in c(1,3,5)) {               # here we call the index j and loop
    print(head(wage1[[j]]))          # through columns of a dataframe
}
@ 
\end{frame}



\begin{frame}{Recommended Literature}

Heiß Chapters 1.9, 3.4, 8.1

Hanck et al. Chapter 5.4
   
\end{frame}
\end{document}


\begin{frame}[fragile]{Example: scatterplot}
\footnotesize
<<cache = F, message= F,out.width='2.45in'>>=
ggplot(wage1) + geom_point(aes(x = educ, y= lwage)) 

@
\end{frame}


\begin{frame}[fragile]{Example: scatterplot}
\texttt{ggplot2} allows to optimize the appearance of the plot
\footnotesize
<<cache = F, message= F,out.width='2.2in'>>=
ggplot(wage1) + geom_point(aes(x = educ, y= lwage), color = "royalblue", alpha = 0.75) + 
  theme_light() + xlab("Years of Education") + ylab("Log Wage")                              + theme(text = element_text(size = 18))
@
\end{frame}


\begin{frame}[fragile]{Example: scatterplot}
choose color of the points based on gender by specifying it within \texttt{aes(...)}
\footnotesize
<<cache = F, message= F,out.width='2.2in'>>=
ggplot(wage1) + geom_point(aes(x = educ, y= lwage, color = factor(female)), alpha = 0.7) + 
  theme_light() + xlab("Years of Education") + ylab("Log Wage")                              + theme(text = element_text(size = 18))
@
\end{frame}

\begin{frame}[fragile]{Example: Adding multiple geom functions}
\small
adding \texttt{geom\_smooth(aes(x = .., y=...), method = "lm")} as an additional geom function, we can add a regression line to the scatter plot
\footnotesize
<<cache = F, message= F,out.width='2.0in'>>=
ggplot(wage1) + geom_point(aes(x = educ, y= lwage), alpha = 0.75) + 
  theme_light() + xlab("Years of Education") + ylab("Log Wage")  + 
  geom_smooth(aes(x = educ, y = lwage), method = "lm", se = F)
  @
\end{frame}

\begin{frame}[fragile]{Example: Adding multiple geom functions}
We can add separate regression lines by gender, by adding it e.g. as color in \texttt{aes(...)}
\vspace{-0.4cm}
\footnotesize
<<cache = F, message= F,out.width='2.2in'>>=
ggplot(wage1) + geom_point(aes(x = educ, y= lwage, color= factor(female)), alpha = 0.75) + 
  theme_light() + xlab("Years of Education") + ylab("Log Wage")  + 
  geom_smooth(aes(x = educ, y = lwage, color = factor(female)), method = "lm", se = F)                        + theme(text = element_text(size = 15))
  @
\end{frame}

\begin{frame}[fragile]{(Some) other geom functions}
\vspace{+0.47cm}


\begin{columns} \small
    \begin{column}{0.53\linewidth}
      \textbf{Density plot}
      \vspace{-0.4cm}
\scriptsize
<<cache = F, message= F,out.width='2.0in'>>=
ggplot(wage1) + 
geom_density(aes(x = wage, color= factor(female))) + 
  theme_light() + xlab("Wage in $1000")   
  @  
\end{column}
\begin{column}{0.47\linewidth}
               \textbf{Bar plot}
      \vspace{-0.4cm}
\footnotesize
<<include= F,out.width='2.2in'>>=
df_means <- as.data.frame(colMeans(wage1)[2:7]) # get means of variables
names(df_means) <- "mean_value"
df_means$variable <- rownames(df_means)
@

<<cache = F, message= F,out.width='2.0in'>>=
ggplot(df_means) + 
geom_col(aes(x= variable, y= mean_value))
@  
\end{column}
\end{columns}
\end{frame}




